\documentclass{article}
\usepackage{amsmath,amssymb,amsthm}
\usepackage{algorithm}
\usepackage[noend]{algpseudocode}
\usepackage{bm}
\usepackage{hyperref}
\usepackage{geometry}
\geometry{margin=1in}
\newtheorem{assumption}{Assumption}
\newtheorem{theorem}{Theorem}
\newtheorem{lemma}{Lemma}
\newcommand{\E}{\mathbb{E}}
\newcommand{\R}{\mathbb{R}}
\newcommand{\norm}[1]{\left\|#1\right\|}
\newcommand{\ip}[2]{\langle #1,#2\rangle}
\newcommand{\grad}{\nabla}
\newcommand{\Hess}{\nabla^{2}}

\begin{document}

\section*{Stochastic Newton Method}

\section*{Mathematical Formulation}
Let $f:\R^{d}\to\R$ be a twice differentiable objective.  
At iteration $t$ we have a stochastic estimate of the Hessian,
\[
\mathbf H_{t}\in\R^{d\times d},
\]
and the (exact) gradient $\grad f(\mathbf x_{t})$.  
With a scalar stepsize $\alpha>0$ the update rule is
\begin{equation}
\boxed{\;
\mathbf x_{t+1}= \mathbf x_{t}-\alpha\,\mathbf H_{t}^{-1}\,\grad f(\mathbf x_{t})\; } \tag{1}
\end{equation}
All expectations below are taken conditioned on the filtration
$\mathcal F_{t}= \sigma(\mathbf x_{0},\mathbf H_{0},\dots,\mathbf H_{t-1})$.

\section*{Pseudocode}
\begin{algorithm}[H]
\caption{Stochastic Newton Method}
\begin{algorithmic}[1]
\State \textbf{Input:} stepsize $\alpha>0$, initial point $\mathbf x_{0}$
\For{$t=0,1,2,\dots$}
    \State Compute (or receive) stochastic Hessian estimate $\mathbf H_{t}$
    \State Compute (or receive) gradient $\mathbf g_{t}= \grad f(\mathbf x_{t})$
    \State $\displaystyle \mathbf x_{t+1}\gets \mathbf x_{t}-\alpha\,\mathbf H_{t}^{-1}\mathbf g_{t}$
\EndFor
\State \textbf{Output:} last iterate $\mathbf x_{T}$ (or average $\frac1T\sum_{t=0}^{T-1}\mathbf x_{t}$)
\end{algorithmic}
\end{algorithm}

\section*{Dry Run Example}
Consider the scalar quadratic
\[
f(w)= (w-3)^{2},\qquad \grad f(w)=2(w-3),\qquad \Hess f(w)=2 .
\]
We use a deterministic ``stochastic’’ Hessian $\mathbf H_{t}=2$ (exact) and stepsize $\alpha=0.1$.
\[
\text{Update: } w_{t+1}= w_{t}-0.1\cdot \frac{1}{2}\,2(w_{t}-3)=w_{t}-0.1(w_{t}-3).
\]

\begin{center}
\begin{tabular}{c|c|c|c}
$t$ & $w_{t}$ & $\grad f(w_{t})$ & $f(w_{t})$\\ \hline
0 & $0$ & $-6$ & $9$\\
1 & $0-0.1\cdot\frac{-6}{2}=0+0.3=0.3$ & $-5.4$ & $(0.3-3)^{2}=7.29$\\
2 & $0.3-0.1\cdot\frac{-5.4}{2}=0.3+0.27=0.57$ & $-4.86$ & $(0.57-3)^{2}=5.9049$\\
3 & $0.57-0.1\cdot\frac{-4.86}{2}=0.57+0.243=0.813$ & $-4.374$ & $(0.813-3)^{2}=4.7777$
\end{tabular}
\end{center}
The iterates move steadily toward the minimiser $w^{\star}=3$ and the objective value decreases.

\newpage
\section*{Assumptions}
\begin{assumption}[Strong convexity and smoothness]\label{ass:strong}
The function $f$ is $\mu$‑strongly convex and $L$‑smooth, i.e.
\[
\forall \mathbf x,\mathbf y\in\R^{d}:\qquad
\frac{\mu}{2}\|\mathbf x-\mathbf y\|^{2}
\le f(\mathbf y)-f(\mathbf x)-\ip{\grad f(\mathbf x)}{\mathbf y-\mathbf x}
\le \frac{L}{2}\|\mathbf y-\mathbf x\|^{2}.
\tag{2}
\]
Consequently $ \mu I \preceq \Hess f(\mathbf x) \preceq L I$ for all $\mathbf x$.
\end{assumption}

\begin{assumption}[Unbiased stochastic Hessian and bounded eigenvalues]\label{ass:hessian}
For every $t$,
\[
\E[\mathbf H_{t}\mid\mathcal F_{t}]=\Hess f(\mathbf x_{t}),
\]
and there exist deterministic scalars $0<\lambda_{\min}\le\lambda_{\max}$ such that
\[
\lambda_{\min} I \preceq \mathbf H_{t}\preceq \lambda_{\max} I
\quad\text{almost surely.}
\tag{3}
\]
Define the condition number of the stochastic Hessian estimator
\[
\kappa\;:=\;\frac{\lambda_{\max}}{\lambda_{\min}}.
\]
\end{assumption}

\begin{assumption}[Exact gradient]\label{ass:grad}
The gradient used in the update is the true gradient,
$\mathbf g_{t}= \grad f(\mathbf x_{t})$.
(If a stochastic gradient is used, an additional unbiasedness and variance bound
must be added; the proof below can be extended in the standard way.)
\end{assumption}

\section*{Main Convergence Theorem}
\begin{theorem}[Linear convergence in expectation]\label{thm:linear}
Under Assumptions \ref{ass:strong}–\ref{ass:grad}, choose a constant stepsize
\[
\alpha\in\bigl(0,\;\tfrac{2\lambda_{\min}}{L}\bigr).
\]
Then for the iterates generated by \eqref{eq:1} we have
\[
\E\!\bigl[\|\mathbf x_{t+1}-\mathbf x^{\star}\|^{2}\mid\mathcal F_{t}\bigr]
\;\le\;
\bigl(1-\rho\bigr)\,\|\mathbf x_{t}-\mathbf x^{\star}\|^{2},
\qquad
\rho:=\frac{2\alpha\mu\lambda_{\min}}{L+\alpha\lambda_{\max}\mu}>0 .
\tag{4}
\]
Consequently,
\[
\E\!\bigl[\|\mathbf x_{t}-\mathbf x^{\star}\|^{2}\bigr]\le (1-\rho)^{t}\,\|\mathbf x_{0}-\mathbf x^{\star}\|^{2},
\qquad
\E\!\bigl[f(\mathbf x_{t})-f(\mathbf x^{\star})\bigr]\le
\frac{L}{2}(1-\rho)^{t}\,\|\mathbf x_{0}-\mathbf x^{\star}\|^{2}.
\tag{5}
\]
\end{theorem}

\section*{Theoretical Properties}
\begin{itemize}
\item \textbf{Stability.} The bound \eqref{eq:4} shows that the method is a contraction
in the mean‑square sense as long as $\alpha$ satisfies the interval above.
\item \textbf{Hyperparameter sensitivity.}
The contraction factor $\rho$ grows linearly with $\alpha$ for small $\alpha$,
but deteriorates when $\alpha$ exceeds $2\lambda_{\min}/L$,
which is the stochastic analogue of the classical Newton step size bound.
\item \textbf{Comparison to SGD.}
For SGD with a diminishing stepsize $O(1/t)$ one obtains a sub‑linear rate $O(1/t)$,
whereas the stochastic Newton method attains a geometric rate $O((1-\rho)^{t})$
provided the Hessian estimator is uniformly well‑conditioned (small $\kappa$).
\end{itemize}

\section*{Discussion}
The theorem hinges on two key ingredients:
\begin{enumerate}
\item Strong convexity guarantees a quadratic lower bound on $f$.
\item Uniform spectral bounds \eqref{eq:3} on the stochastic Hessian keep the
inverse $\mathbf H_{t}^{-1}$ from blowing up and allow us to relate the
Newton direction to the exact Newton direction.
\end{enumerate}
If the estimator is biased or its eigenvalues are only bounded in
expectation, the proof can be adapted by inserting additional
martingale concentration arguments (cf. standard stochastic Newton
analyses).

\newpage
\section*{Preliminaries (Definitions and Lemmas)}
\begin{lemma}[Quadratic growth]\label{lem:quad}
Under Assumption \ref{ass:strong},
\[
\frac{\mu}{2}\|\mathbf x-\mathbf x^{\star}\|^{2}\le f(\mathbf x)-f(\mathbf x^{\star})
\le\frac{L}{2}\|\mathbf x-\mathbf x^{\star}\|^{2}.
\tag{6}
\]
\end{lemma}
\begin{proof}
Directly from the definition of $\mu$‑strong convexity and $L$‑smoothness
by setting $\mathbf y=\mathbf x^{\star}$ and using $\grad f(\mathbf x^{\star})=0$.
\end{proof}

\begin{lemma}[Spectral bounds on the inverse]\label{lem:inv}
If $\lambda_{\min}I\preceq\mathbf H\preceq\lambda_{\max}I$, then
\[
\frac{1}{\lambda_{\max}}I\preceq\mathbf H^{-1}\preceq\frac{1}{\lambda_{\min}}I.
\tag{7}
\]
\end{lemma}
\begin{proof}
Take any $\mathbf v\neq0$ and write
\[
\lambda_{\min}\|\mathbf v\|^{2}\le\ip{\mathbf H\mathbf v}{\mathbf v}\le\lambda_{\max}\|\mathbf v\|^{2}.
\]
Dividing by $\ip{\mathbf H\mathbf v}{\mathbf v}$ and using $\mathbf H^{-1}$ yields
\[
\frac{1}{\lambda_{\max}}\|\mathbf v\|^{2}\le\ip{\mathbf H^{-1}\mathbf v}{\mathbf v}\le
\frac{1}{\lambda_{\min}}\|\mathbf v\|^{2}.
\]
Since the inequality holds for all $\mathbf v$, the matrix inequality \eqref{eq:7}
follows.\qedhere
\end{proof}

\section*{Main Proof (Step‑by‑Step Derivation)}
We prove Theorem \ref{thm:linear}.  
All expectations are conditional on $\mathcal F_{t}$ unless otherwise stated.

\begin{enumerate}
\item[Step 1.]  Write the error vector $\mathbf e_{t}:=\mathbf x_{t}-\mathbf x^{\star}$.  
Using the update \eqref{eq:1},
\[
\mathbf e_{t+1}= \mathbf e_{t}-\alpha\mathbf H_{t}^{-1}\grad f(\mathbf x_{t}).
\tag{8}
\]

\item[Step 2.]  Square the norm and expand:
\[
\begin{aligned}
\|\mathbf e_{t+1}\|^{2}
&= \|\mathbf e_{t}\|^{2}
-2\alpha\ip{\mathbf H_{t}^{-1}\grad f(\mathbf x_{t})}{\mathbf e_{t}}
+\alpha^{2}\|\mathbf H_{t}^{-1}\grad f(\mathbf x_{t})\|^{2}.
\end{aligned}
\tag{9}
\]

\item[Step 3.]  Replace $\grad f(\mathbf x_{t})$ by the integral form of the gradient
(using the fundamental theorem of calculus for vector‑valued functions):
\[
\grad f(\mathbf x_{t})
= \bigl(\int_{0}^{1}\Hess f(\mathbf x^{\star}+s\mathbf e_{t})\,\mathrm ds\bigr)\mathbf e_{t}
=: \mathbf A_{t}\mathbf e_{t},
\tag{10}
\]
where $\mathbf A_{t}$ is a symmetric positive‑definite matrix satisfying
\[
\mu I\preceq\mathbf A_{t}\preceq L I
\quad\text{(by Assumption \ref{ass:strong})}.
\tag{11}
\]

\item[Step 4.]  Substitute \eqref{eq:10} into \eqref{eq:9}:
\[
\|\mathbf e_{t+1}\|^{2}
= \|\mathbf e_{t}\|^{2}
-2\alpha\ip{\mathbf H_{t}^{-1}\mathbf A_{t}\mathbf e_{t}}{\mathbf e_{t}}
+\alpha^{2}\|\mathbf H_{t}^{-1}\mathbf A_{t}\mathbf e_{t}\|^{2}.
\tag{12}
\]

\item[Step 5.]  Apply Lemma \ref{lem:inv} and the spectral bounds \eqref{eq:11}:
\[
\begin{aligned}
\ip{\mathbf H_{t}^{-1}\mathbf A_{t}\mathbf e_{t}}{\mathbf e_{t}}
&= \ip{\mathbf A_{t}^{1/2}\mathbf H_{t}^{-1}\mathbf A_{t}^{1/2}\mathbf A_{t}^{1/2}\mathbf e_{t}}
      {\mathbf A_{t}^{-1/2}\mathbf e_{t}}  \\
&\ge \frac{1}{\lambda_{\max}}\,
      \ip{\mathbf A_{t}\mathbf e_{t}}{\mathbf e_{t}}
\ge \frac{\mu}{\lambda_{\max}}\|\mathbf e_{t}\|^{2}.
\end{aligned}
\tag{13}
\]
Similarly,
\[
\begin{aligned}
\|\mathbf H_{t}^{-1}\mathbf A_{t}\mathbf e_{t}\|^{2}
&= \ip{\mathbf H_{t}^{-1}\mathbf A_{t}\mathbf e_{t}}
        {\mathbf H_{t}^{-1}\mathbf A_{t}\mathbf e_{t}}   \\
&\le \frac{1}{\lambda_{\min}^{2}}\,
      \|\mathbf A_{t}\mathbf e_{t}\|^{2}
\le \frac{L^{2}}{\lambda_{\min}^{2}}\|\mathbf e_{t}\|^{2}.
\end{aligned}
\tag{14}
\]

\item[Step 6.]  Plug the bounds \eqref{eq:13}–\eqref{eq:14} into \eqref{eq:12}:
\[
\|\mathbf e_{t+1}\|^{2}
\le \Bigl(1-2\alpha\frac{\mu}{\lambda_{\max}}
      +\alpha^{2}\frac{L^{2}}{\lambda_{\min}^{2}}\Bigr)\|\mathbf e_{t}\|^{2}.
\tag{15}
\]

\item[Step 7.]  Choose $\alpha$ such that the coefficient in \eqref{eq:15}
is strictly smaller than one.  The quadratic in $\alpha$ is minimized at
$\alpha^{\star}= \frac{\mu\lambda_{\min}^{2}}{L^{2}\lambda_{\max}}$.
For any $\alpha\in\bigl(0,\frac{2\lambda_{\min}}{L}\bigr)$ we have
\[
1-2\alpha\frac{\mu}{\lambda_{\max}}
      +\alpha^{2}\frac{L^{2}}{\lambda_{\min}^{2}}
= 1-\rho,
\qquad
\rho:=\frac{2\alpha\mu\lambda_{\min}}{L+\alpha\lambda_{\max}\mu}>0 .
\tag{16}
\]
Thus,
\[
\E[\|\mathbf e_{t+1}\|^{2}\mid\mathcal F_{t}]
\le (1-\rho)\,\|\mathbf e_{t}\|^{2}.
\tag{17}
\]

\item[Step 8.]  Take total expectation and iterate the contraction:
\[
\E\|\mathbf e_{t}\|^{2}\le (1-\rho)^{t}\,\|\mathbf e_{0}\|^{2}.
\tag{18}
\]

\item[Step 9.]  Convert the distance bound into an objective‑value bound
using Lemma \ref{lem:quad} (upper part):
\[
f(\mathbf x_{t})-f(\mathbf x^{\star})
\le \frac{L}{2}\|\mathbf e_{t}\|^{2}
\le \frac{L}{2}(1-\rho)^{t}\|\mathbf e_{0}\|^{2}.
\tag{19}
\]
This completes the proof of Theorem \ref{thm:linear}. \qed
\end{enumerate}

\section*{Rate Derivation (With Constants)}
From \eqref{eq:16} the linear rate constant can be written explicitly as
\[
\rho(\alpha)=\frac{2\alpha\mu\lambda_{\min}}
                  {L+\alpha\lambda_{\max}\mu}
          =\frac{2\alpha\mu}{L}\cdot
            \frac{\lambda_{\min}}{\lambda_{\min}+ \alpha\mu\kappa^{-1}},
\qquad\kappa=\frac{\lambda_{\max}}{\lambda_{\min}}.
\]
If we set the stepsize to the maximal admissible value
$\alpha= \frac{2\lambda_{\min}}{L}$,
\[
\rho_{\max}= \frac{4\mu\lambda_{\min}^{2}}
                  {L^{2}+2\mu\lambda_{\min}\lambda_{\max}}
            = \frac{4\mu}{L}\cdot\frac{1}{\kappa+2\mu/L}.
\]
Hence a well‑conditioned Hessian estimator (small $\kappa$) yields a
contraction factor close to the deterministic Newton rate $1-\frac{2\mu}{L}$.

\section*{Special Cases}
\begin{itemize}
\item \textbf{Exact Newton method.}
If $\mathbf H_{t}=\Hess f(\mathbf x_{t})$ (so $\lambda_{\min}=\mu$, $\lambda_{\max}=L$)
and $\alpha=1$, then $\rho= \frac{2\mu L}{L^{2}+ \mu L}= \frac{2\mu}{L+\mu}$,
recovering the classic linear rate for Newton’s method on strongly convex
quadratics.
\item \textbf{Identity preconditioner.}
When $\mathbf H_{t}=I$, we have $\lambda_{\min}=\lambda_{\max}=1$,
$\kappa=1$, and \eqref{eq:16} reduces to the standard gradient descent
contraction $1-2\alpha\mu+ \alpha^{2}L^{2}$, reproducing the well‑known
condition $\alpha\in(0,2/L)$.
\end{itemize}

\end{document}